\documentclass[10pt]{article}
\usepackage[T1]{fontenc}
\usepackage[utf8]{inputenc}
\usepackage[french]{babel}
\usepackage{amsmath,amssymb,mathrsfs,mathpazo}
\usepackage{geometry}
\geometry{a4paper,margin=2cm}
\usepackage{graphicx,array,tabularx,tabularray,colortbl,multirow,multicol,booktabs,enumitem,pifont,ulem,fancybox,diagbox,cancel}
\usepackage{pstricks,pstricks-add,pst-plot,pst-tree,pst-node,pst-text,pst-func,pst-math,pst-bspline,pst-3dplot,pst-eucl,pst-solides3d}
\let\pscalculate\undefined
\usepackage{pst-bezier}
\usepackage{tikz,pgfplots,tkz-tab}
\usetikzlibrary{arrows,arrows.meta,patterns,calc,positioning,trees,decorations.markings,shapes.geometric,shapes.misc,shapes.arrows,intersections,decorations.pathmorphing,decorations.pathreplacing,backgrounds,shadows,babel}
\pgfplotsset{compat=1.18}
\usepgfplotslibrary{fillbetween}
\usepackage[np]{numprint}
\usepackage{listings,makecell,pas-tableur,twoopt}
\let\barre\undefined
\usepackage[pstricks]{tablvar}
\definecolor{darkgreen}{RGB}{0,100,0}
\providecommand{\red}{\color{red}}
\providecommand{\blue}{\color{blue}}
\providecommand{\green}{\color{green}}
\providecommand{\black}{\color{black}}
\providecommand{\white}{\color{white}}
\providecommand{\gray}{\color{gray}}
\providecommand{\cyan}{\color{cyan}}
\providecommand{\magenta}{\color{magenta}}
\providecommand{\yellow}{\color{yellow}}
\providecommand{\violet}{\color{violet}}
\providecommand{\orange}{\color{orange}}
\providecommand{\euro}{\texteuro}
\providecommand{\EUR}{\texteuro}
\providecommand{\ieme}{\textsuperscript{e}}
\providecommand{\degres}{\ensuremath{^\circ}}
\providecommand{\degree}{\ensuremath{^\circ}}
\providecommand{\vv}[1]{\overrightarrow{#1}}
\providecommand{\Coord}[1]{#1}
\providecommand{\decosix}{}
\providecommand{\starredbullet}{\ensuremath{\bullet}}
\pagestyle{empty}
\setlength{\parindent}{0pt}
\usepackage{adjustbox}\DeclareUnicodeCharacter{2236}{:}\DeclareUnicodeCharacter{2190}{\ensuremath{\leftarrow}}
\providecommand{\R}{\mathbb{R}}
\providecommand{\N}{\mathbb{N}}
\providecommand{\Z}{\mathbb{Z}}
\providecommand{\Q}{\mathbb{Q}}
\providecommand{\D}{\mathbb{D}}
\providecommand{\C}{\mathbb{C}}
\providecommand{\e}{{\mathrm{e}}}
\providecommand{\vect}[1]{\overrightarrow{#1}}
\providecommand{\vectt}[1]{\overrightarrow{\mathrm{#1}}}
\providecommand{\vvt}[1]{\overrightarrow{\mathrm{#1}}}
\providecommand{\np}[2][]{{#2~#1}}
\providecommand{\vv}[1]{{\overrightarrow{#1}}}
\renewcommand{\d}{{\mathrm{d}}}
\providecommand{\degres}{{^\circ}}
\providecommand{\degre}{{^\circ}}
\providecommand{\blue}{\color{blue}}
\providecommand{\red}{\color{red}}
\providecommand{\enskip}{\quad}
\providecommand{\llbracket}{\lbrack\!\lbrack}
\providecommand{\rrbracket}{\rbrack\!\rbrack}
\providecommand{\euro}{{\text{€}}}
\providecommand{\no}{{\text{n°}}}
\providecommand{\interval}[2]{{[#1;#2]}}
\providecommand{\intervFO}[2]{{[#1;#2[}}
\providecommand{\dotfill}{\ldots}

\providecommand{\ds}{\displaystyle}
\providecommand{\textding}[1]{#1}
\providecommand{\Oij}{\(\mathrm{O};\vec{\imath},\vec{\jmath}\)}
\providecommand{\Oijk}{\(\mathrm{O};\vec{\imath},\vec{\jmath},\vec{k}\)}
\providecommand{\Ouv}{\(\mathrm{O};\vec{u},\vec{v}\)}
\providecommand{\intervFF}[2]{\ensuremath{[#1;#2]}}
\renewcommand{\euro}{\ifmmode\text{\texteuro}\else\texteuro\fi}
\DeclareUnicodeCharacter{00B0}{\ifmmode^{\circ}\else\textdegree\fi}
\DeclareUnicodeCharacter{20AC}{\ifmmode\text{\texteuro}\else\texteuro\fi}
% Énoncés : Éducation nationale. Transcription et corrigés : APMEP, Denis Vergès et contributeurs. Découpage et classement : CoopMaths, Sébastien Lozano, Éric Elter et contributeurs. Adaptation de la présentation : Maths Explorer.
% CC BY-SA : https://coopmaths.fr/www/about/
% https://coopmaths.fr/alea/static/bac/2021/tex/bac_2021_06_etrangers_3.tex
\begin{document}
\noindent\textit{APMEP / CoopMaths. CC BY-SA. Présentation adaptée par Maths Explorer.}\par\medskip
\textbf{Commun à tous les candidats}

\medskip

En mai 2020, une entreprise fait le choix de développer le télétravail afin de s'inscrire dans une démarche écoresponsable.

Elle propose alors à ses \np{5000}~collaborateurs en France de choisir entre le télétravail et le travail au sein des locaux de l'entreprise.

En mai 2020, seuls $200$ d'entre eux ont choisi le télétravail.

Chaque mois, depuis la mise en place de cette mesure, les dirigeants de l'entreprise constatent que $85$\,\% de ceux qui avaient choisi le télétravail le mois précédent choisissent de continuer, et que, chaque mois, $450$ collaborateurs supplémentaires choisissent le télétravail.

On modélise le nombre de collaborateurs de cette entreprise en télétravail par la suite $\left(a_n\right)$.

Le terme $a_n$ désigne ainsi une estimation du nombre de collaborateurs en télétravail le $n$-ième mois après le mois de mai 2020. Ainsi $a_0 = 200$.

\bigskip

\textbf{Partie A }:

\medskip

\begin{enumerate}
\item Calculer $a_1$.
\item Justifier que pour tout entier naturel $n$,\,
$a_{n+1} = 0,85a_n + 450$.
\item On considère la suite $\left(v_n\right)$ définie pour tout entier naturel n par: $v_n = a_n - \np{3000}$.
	\begin{enumerate}
		\item Démontrer que la suite $\left(v_n\right)$ est une suite géométrique de raison $0,85$. 
		\item Exprimer $v_n$ en fonction de $n$ pour tout entier naturel $n$.
		\item En déduire que, pour tout entier naturel $n$,\,
$a_n = \np{- 2800} \times  0,85^n + \np{3000}$.
	\end{enumerate}
\item Déterminer le nombre de mois au bout duquel le nombre de télétravailleurs sera strictement supérieur à \np{2500}, après la mise en place de cette mesure dans l'entreprise.
\end{enumerate}

\bigskip

\textbf{Partie B }:

\medskip

Afin d'évaluer l'impact de cette mesure sur son personnel, les dirigeants de l'entreprise sont parvenus à modéliser le nombre de collaborateurs satisfaits par ce dispositif à l'aide de la suite 
$\left(u_n\right)$ définie par $u_0 = 1$ et, pour tout entier naturel $n$,

\[u_{n+1}  = \dfrac{5u_n + 4}{u_n + 2}\]

où $u_n$ désigne le nombre de milliers de collaborateurs satisfaits par cette nouvelle mesure au bout de $n$ mois après le mois de mai 2020.

\medskip

\begin{enumerate}
\item Démontrer que la fonction $f$ définie pour tout $x \in  [0~;~ +\infty[$ par $f(x) = \dfrac{5x+4}{x+2}$ est 
strictement croissante sur $[0~;~+ \infty[$.
\item 
	\begin{enumerate}
		\item Démontrer par récurrence que pour tout entier naturel 
		\[0 \leqslant u_n \leqslant u_{n+1}  \leqslant  4.\]
		
		\item Justifier que la suite $\left(u_n\right)$ est convergente.
	\end{enumerate}
\item On admet que pour tout entier naturel $n$,

\[0  \leqslant 4 - u_n \leqslant 3 \times \left(\dfrac{1}{2}\right)^n.\]

En déduire la limite de la suite $\left(u_n\right)$ et l'interpréter dans le contexte de la modélisation.
\end{enumerate}

\bigskip

\end{document}
